\section{Polygonal knots in three and four dimensions}
\label{sec:geometry}
A polygonal link is a list of cyclic lists of rational points. Each circuit has at least three edges. Adjacent edges meet precisely at their common endpoint; all other edges are disjoint. In dimension $d$, an elementary triangle move replaces one side of a nondegenerate rational triangle by the other two or vice versa, when the triangle meets the rest of the polygon only in the permitted endpoints. Subdivision and its inverse are allowed. These are assertions about finite lists. Intersections of rational simplices reduce to feasibility of rational linear equations and inequalities in barycentric coordinates. For a fixed ambient dimension and simplex size, eliminate one variable at a time by combining every lower bound with every upper bound; retain the open/closed flags if a facet is excluded. Given values for the remaining variables, back-substitute a midpoint between bounds, the common value of two closed equal bounds, or a point one unit beyond a one-sided bound. This proves both the test and rational witness property. In particular disjointness is a decidable arithmetic predicate.

\begin{lemma}[Finite grid avoidance]\label{lem:grid}
A nonzero rational polynomial of total degree at most $d$ in $k$ variables has a nonzero value on $\{0,\ldots,d\}^k$.
\end{lemma}
\begin{proof}
In one variable, division by $X-a$ with remainder $f(a)$ shows that each different root extracts a new linear factor. A polynomial having $d+1$ distinct roots and degree at most $d$ is zero. In several variables regard the last variable as the main variable. Vanishing at the entire grid forces all coefficient polynomials to vanish at every grid point in the other variables. Repeat. For its arithmetic formulation inspect the bounded grid for a witness. If none is found the coefficient argument gives a contradiction. For a varying number of variables, fix the input coefficient array and aggregate the tests on all its successive coefficient slices in a bounded conjunction. There are at most $(d+1)^k$ coefficient slots at any level. Induct on the number of coordinates eliminated with this conjunction as invariant. The projection application uses only two variables. For products we use the lexicographically largest nonzero coefficient of each factor: their product is the largest coefficient of the product. All these assertions concern finite coefficient arrays.
\end{proof}

\begin{theorem}[Regular rational projection, with a bound]\label{th:projection}
From an embedded rational polygonal link with $N$ edges in three-space one can compute $u,v\in\{0,\ldots,10(N+1)^3\}$ such that
\[
             (x,y,z)\longmapsto (x-uz,y-vz)
\]
is a regular diagram projection. Overpass bits are recovered by rational comparisons. The crossing count is at most $\binom{N}{2}$, and at most $N(N-3)/2$ for a knot.
\end{theorem}
\begin{proof}
Write $d=(u,v,1)$. We describe finitely many nonzero polynomial factors to be avoided. For a nonzero edge direction forbid collapse by taking a nonzero component polynomial of its projected direction. Do the same for differences of distinct vertices. For each nonparallel pair of supporting lines with directions $a,b$ exclude $d\cdot(a\times b)=0$. For a pair of distinct parallel lines exclude directions in their common plane of directions: its nonzero normal is $a\times(Q-P)$, where $P,Q$ are on the two lines. Collinear segments retain their separation and order whenever the common line does not collapse. For a vertex outside a supporting line forbid the plane of directions spanned by the edge and the vector from a point of the line to the vertex. A vertex on the supporting line is already handled by injectivity there. All plane factors have degree at most one and a nonzero normal gives a polynomial which is not identically zero.

Two different nonparallel coplanar lines meet in three-space. If their projections meet in segment interiors, their unique projected intersection must be the image of that very point. Since the original segments are embedded, this is impossible. Parallel pairs were handled above. Consequently a triple crossing of segment interiors can involve only three pairwise skew supporting lines. For such a triple, write the three projected line equations in homogeneous plane coordinates, each obtained as the cross product of its two projected endpoints $(x-uz,y-vz,1)$. Their three-row determinant $H(u,v)$ detects concurrence. Each row has entries affine in $(u,v)$ because the two $uv$ terms cancel, so $\deg H\le3$.

Here is why this last factor really is nonzero. Take a rational point $P$ on the first line and let $Q$ vary on the second. The plane through $P$ and line two cannot contain line three, since lines two and three are skew. In its projective completion it therefore meets line three in a single point $S$ (possibly a direction at infinity). This is just a linear equation along line three; its constant nonzero case gives its point at infinity. $S\ne P$. At most one of the lines $PQ$ passes through $S$. Choose $P$ so that the whole second line is not at its height when that line is horizontal. This is possible since two horizontal lines at the same height cannot be skew. Then at most one further $Q$ gives equal heights. Among three distinct rational $Q$'s choose one avoiding both exceptions. In the direction $Q-P$, normalized to last coordinate one, the projections of the first two lines meet transversely: coplanarity of their directions with $Q-P$ would say the lines were coplanar. The third does not collapse to a point, since a direction parallel to it would put its point at infinity on $PQ$ and in the chosen plane. The fiber $PQ$ misses line three. Hence its projected line misses the finite intersection of the first two and $H(u,v)\ne0$. This witness construction uses only affine linear equations (at most two candidate points suffice for $P$).

The sum of the degrees of these factors is bounded by $10(N+1)^3$: there are at most $N$ edges and vertices, quadratically many point/line and pair conditions, and at most $\binom N3$ determinants of degree three. Apply Lemma~\ref{lem:grid} to their product. Distinct projected intersections on one edge can now be ordered by rational edge parameter. At a crossing the original heights cannot coincide, by injectivity of the original link. Only nonadjacent edges can cross, which gives the two numerical bounds. The linear shear $(x-suz,y-svz,z)$ has an explicit inverse at every $s\in[0,1]$, giving the conventional ambient interpretation of our coordinate choice.
\end{proof}

\begin{proposition}[Lifting a diagram]
Every rational polygonal diagram in the plane has a rational polygonal link lifting it. Lifts with the same parametrized projection and the same strict height order at each crossing can be joined by a linear interpolation of heights through embeddings.
\end{proposition}
\begin{proof}
Subdivide each original edge at every crossing occurrence. Assign height $+1$ to the upper occurrence, $-1$ to the lower and $0$ to the remaining vertices. Interpolate along the subdivided edges. A coincidence in space would project to a coincidence in the plane; at the only distinct preimages of any point the heights differ. For two lifts, take a common subdivision. If the two differences of upper and lower heights are $a,b>0$, the interpolated difference $(1-s)a+sb$ is at least $\min(a,b)>0$. This proves the assertion with rational input even when reading the interpolation over real parameters. The arithmetic certificate consists of the finite list of endpoint inequalities and affine formulas.
\end{proof}

\begin{theorem}[Unknotting every polygon in four-space]\label{th:four}
A rational polygonal knot in $\mathbb Q^4$, with $N$ edges, can be changed to a triangle using at most $N-1$ elementary triangle moves. A cone point with coordinates $(t,t^2,t^3,t^4)$ exists with $0\le t\le4N^2$.
\end{theorem}
\begin{proof}
For every pair of edges, including a repeated edge, their four endpoints span an affine space of dimension at most three. By rational row reduction choose a nonzero affine linear form vanishing on that space. Its value on $(t,t^2,t^3,t^4)$ is a nonzero polynomial of degree at most four, as the five monomials are distinct. The product for at most $N^2$ forms has at most $4N^2$ distinct roots. Bounded search produces $P$ avoiding all these affine spans. Cone every edge to $P$.

Two points $x,y$ on the old polygon lying on a line through $P$ must coincide: otherwise $P$ belongs to the affine span of the edges containing them. Thus two cone triangles intersect only in their prescribed common faces, as follows by inspecting rays from $P$. Each triangle is nondegenerate. Write the vertices $v_0,\ldots,v_{N-1}$. Replace $v_0v_1$ by $v_0Pv_1$. Next replace $Pv_1v_2$ by $Pv_2$, then $Pv_2v_3$ by $Pv_3$, and continue until the boundary is the triangle $v_0Pv_{N-1}v_0$. These moves are clean by the intersection statement. There are $1+(N-2)=N-1$ moves (or keep an initial triangle). Any nondegenerate rational triangle is a standard unknotted polygon: inside its plane it bounds its filled affine disk. If one fixed standard triangle is desired, replace its vertices one at a time via auxiliary triangles as detailed below.
\end{proof}

\begin{center}
\begin{tikzpicture}[scale=.79,line join=round]
\coordinate (o) at (0,0);
\foreach \i/\a in {0/90,1/18,2/-54,3/-126,4/162}
  \coordinate (v\i) at (\a:1.55);
\fill[sea!15] (v0)--(o)--(v1)--cycle;
\foreach \i in {0,...,4} \draw[sea!70] (o)--(v\i);
\draw[ink,line width=1.1pt] (v0)--(v1)--(v2)--(v3)--(v4)--cycle;
\node[fill=white,inner sep=0.9pt] at (o) {$P$};
\foreach \i/\a in {0/90,1/18,2/-54,3/-126,4/162}
  \node at (\a:1.90) {$v_{\i}$};
\end{tikzpicture}
\end{center}
\noindent The pentagon shows the incidence pattern of the cone disk. In the shaded triangle, replace the boundary edge $v_0v_1$ by $v_0Pv_1$. Each subsequent move deletes one old vertex along an adjacent triangle. Avoidance of the edge-pair spans proves that these faces are embedded in four-space.

For this last identification there is also a direct ambient certificate. Complete the two independent side vectors to a rational basis, reversing the sign of a complementary vector if necessary to make its determinant positive. Translation followed by this change of basis takes a fixed coordinate triangle onto the given triangle. Gaussian elimination expresses a positive-determinant rational matrix as shears, positive diagonal scalings, and simultaneous sign changes on pairs of axes. A signed row exchange and a pair sign change are products of shears (in two dimensions $\left(\begin{smallmatrix}0&1\\-1&0\end{smallmatrix}\right)$ is the product of shears with successive off-diagonal entries $1,-1,1$). Each shear and positive diagonal scaling interpolates to the identity rationally without losing invertibility. For a triangle-move-only version, join the original and target triangles via a third triangle with new rational vertices chosen successively off the affine spans of every triple of already named points and off their lines. Use the moment curve test above with the degree bound four per forbidden plane. Replace vertices one at a time, first to the auxiliary triple and then to the target triple. A replacement $ABC\to A'BC$ uses two faces of the tetrahedron $AA'BC$: insert $A'$ on $AB$ by a triangle move, then delete $A$ along $CA A'$. Affine independence of every mixed quadruple used here follows from the choices of the auxiliary points (take all six old and target vertices as named at the outset; the old and target triples themselves are noncollinear). This gives an explicit connection to a fixed triangle.

\paragraph{Internal bounds.}
Solving the fixed small linear systems above uses determinants of bounded size, with signs and zero tests on integers. For the regular projection one may simply try every pair in the stated square and validate the resulting diagram. For the cone point there are $4N^2+1$ trials. To obtain the triangle-move certificate, write the remaining prefix fan after each edge: its validity is the bounded conjunction of the pairwise simplex tests. Arithmetic induction is on the number of edges already removed, with this Boolean conjunction and the explicit boundary list as invariant. No compactness or topological general-position theorem enters either construction.

\subsection{A ribbon disk gives a slice disk}
An oriented polygonal knot is \emph{PL slice}, in the sense used here, if it bounds a properly embedded rational triangulated disk in a four-dimensional box, with its boundary in the interior of a three-dimensional face. A \emph{ribbon disk certificate} is a triangulated parameter disk and a rational simplexwise affine immersion $f$ into three-space with the following recorded self-intersections. The double locus has finitely many disjoint arcs and no triple points. Each double arc has two disjoint preimage arcs mapped homeomorphically onto it, one proper (ends on the boundary) and the other entirely in the interior. These arcs are subcomplexes after a supplied subdivision, and the lists account for every coincidence of distinct points, as can be checked by triangle-pair intersections. The boundary is the prescribed knot.

\begin{theorem}[Ribbon implies slice, with an explicit fourth coordinate]
From a ribbon disk certificate one can construct a slice disk certificate by one barycentric subdivision and a vertex-height assignment.
\end{theorem}
\begin{proof}
Barycentrically subdivide the parameter disk, retaining its boundary and all the marked arcs as subcomplexes. Give height $0$ to boundary vertices, $2$ to vertices of every interior preimage arc, and $1$ to every other vertex. Extend height $h$ affinely on each small triangle. Its zero set is exactly the boundary: vertices of a barycentric simplex form a chain of nested old faces, and any face of the simplex all of whose vertices are boundary barycenters is contained in the largest boundary face in that chain. On each interior preimage arc $h=2$ everywhere. On its paired proper arc $0\le h\le1$, including at its endpoints. Thus $(f,h)$ separates every former double point, and its projection precludes any new coincidences. It gives a rational embedded triangulated disk; its inverse on each image simplex is the affine inverse furnished by the nondegenerate input triangle. Pick an integer $M$ strictly larger than every absolute coordinate of every image vertex. The disk lies in $[-M,M]^3\times[0,3]$, meets the boundary of this box exactly in the given knot at height zero, and has interior in the open box. A triangle yields six barycentric triangles, and all checks are bounded vertex, face and intersection tests. This proves the forward implication for the finite PL definitions; the converse is the slice--ribbon question, outside these proofs.
\end{proof}
